
\begin{ledgroupsized}[r]{120mm}
\footnotesize
\pstart
\noindent\textbf{Überlieferung:}
\pend
\end{ledgroupsized}
\begin{ledgroupsized}[r]{114mm}
\footnotesize
\pstart
\parindent -6mm
\makebox[6mm][l]{\textit{L}}%
Konzept:
LH~XXXVII~5 Bl.~21. 
Ein Blatt~8\textsuperscript{o};
Papiererhaltungsmaßnahmen;
alle Ränder beschnitten.
Eindreiviertel Seiten.
\pend
\end{ledgroupsized}
\selectlanguage{latin}
\vspace{8mm}
\pstart
\noindent
\href{\Scan{LH_37_05_021r}}%
{\lbrack21~r\textsuperscript{o}\rbrack}
\hfill
Martii 1678. 
\pend
\nobreak
\pstart
\centering De centro gravitatis.\indexS{centrum gravitatis}
\pend
\nobreak
\vspace{0.5\baselineskip}
\pstart
\noindent 
Quod in quavis linea 
\edtext{sit punctum}{\lemma{sit}\Bfootnote{%
\textit{(1)}~centru %
\textit{(2)}~punctum~\textit{L}}}
aequilibrii,\indexS{punctum aequilibrii} 
in quavis vero superficie axis aequilibrii\indexS{axis aequilibrii} facile patet. 
Et ad principium vectis\indexS{principium vectis} 
demonstrandum\indexS{demonstratio principii vectis} non puto
opus esse consideratione centri gravitatis\indexS{centrum gravitatis} in plano vel solido, 
in quo sese axes 
\edtext{aequilibrii\indexS{axis aequilibrii} intersecant.\indexS{intersectio axium aequilibrii}}{\lemma{aequilibrii}\Bfootnote{%
\textit{(1)}~, cum %
\textit{(2)}~intersecant.~\textit{L}}}
\pend
\begin{edfigure}[1em][1em]
\includegraphics[scale=1]{%
gesamttex/edit_VIII-4/images/LH_37_05_021_d1_021r.pdf}
\edcaption{\lbrack\textit{Fig.~1}\rbrack}
\end{edfigure}
\pstart
Si sint duo axes 
\edtext{aequilibrii\indexS{axis aequilibrii} \textit{AB}, \textit{CD} sese ad angulos rectos secantes in \textit{L} ac ducatur \textit{ELF},}{%
\lemma{aequilibrii}%
\Bfootnote{%
\textit{(1)}~sese ad angulos rectos secantes \textit{AB}, \textit{CD}
\textit{(2)}~\textit{AB}, \textit{CD} \lbrack...\rbrack\ ducatur 
\textit{(a)}~\textit{EF} %
\textit{(b)}~\textit{ELF} %
\textit{(c)}~\textit{ELF},~\textit{L}}}
ita ut angulus \textit{ELC} sit semirectus\lbrack,\rbrack\ 
ajo \textit{EF} fore etiam axem aequilibrii,\indexS{axis aequilibrii} 
alioqui enim corpus si inclinet in latus \textit{CB} verbi gratia tunc magis inclinabit in \textit{C} quam in \textit{D}, et in \textit{B} quam in \textit{A}. 
Idem est de \textit{GH}. Et ita porro pergendo demonstrari potest si rursus bisecetur angulus \textit{ALE}, per 
\edtext{rectam \textit{MLN}}{%
\lemma{rectam}%
\Bfootnote{%
\textit{(1)}~\textit{MN}
\textit{(2)}~\textit{MLN}~\textit{L}}}
fore etiam \textit{MN} axem aequilibrii,\indexS{axis aequilibrii} nam si inclinet in partem \textit{ECG} 
\edtext{utique magis}{\lemma{utique}\Bfootnote{%
\textit{(1)}~non ratio %
\textit{(2)}~ubique %
\textit{(3)}~\bars{utique \textit{streicht Hrsg.}} magis~\textit{L}}}
inclinabit in \textit{E} quam in \textit{F} et in \textit{C} quam in \textit{D}, et in \textit{G} quam in \textit{H}, et in \textit{B} quam in \textit{A} 
et ita continuata bisectione in infinitum, ostendetur quamlibet rectam per \textit{L} transeuntem esse axem aequilibrii,\indexS{axis aequilibrii} 
aut ab eo minus differre angulo quolibet assignabili. 
Hanc ratiocinationem pro demonstratione habeo, nam parem ubique rationem esse manifestum est. 
Imo ne opus est quidem consideratione progressus in infinitum,\indexS{progressus in infinitum} 
si initio non sint recti anguli, sed quicunque dati postea bisecentur.
\pend
\pstart
Idem videtur demonstrari 
\edtext{posse, magis}{\lemma{posse,}\Bfootnote{%
\textit{(1)}~si %
\textit{(2)}~magis~\textit{L}}}
geometrice minusque metaphysice, si componamus planum ex multis lineis\indexS{lineae sese secantes in centro gravitatis}
sese secantibus in puncto \textit{L} (quibus lineis aliquam crassitiem dabimus\lbrack)\rbrack, et centra gravitatis singularum\indexS{centrum gravitatis singulorum}
\href{\Scan{LH_37_05_021v}}%
{\lbrack21~v\textsuperscript{o}\rbrack}
quaerendo et inter se connectendo ostendamus tandem, unum omnium fore centrum gravitatis,\indexS{centrum gravitatis omnium} 
sed ad hanc rem Lemmate prius opus est, quod centrum gravitatis centrorum gravitatis\indexS{centrum gravitatis centrorum gravitatis} est centrum gravitatis totius.\indexS{centrum gravitatis totius}
\pend
\pstart
Corpus quod ab omni latere uniforme est, ut Sphaera,\indexS{Sphaera} utique habet centrum gravitatis,\indexS{centrum gravitatis} jam duo corpora centrum gravitatis habentia, utique habent centrum gravitatis 
\edtext{commune,\indexS{centrum gravitatis commune} alicubi}{\lemma{commune,}\Bfootnote{%
\textit{(1)}~in %
\textit{(2)}~alicubi~\textit{L}}}
in rigida recta earum centra conjungente.\indexS{recta centra gravitatis conjungens} 
Et ita invento eorum centro gravitatis 
\edtext{communi,\indexS{centrum gravitatis commune} id}{\lemma{communi,}\Bfootnote{%
\textit{(1)}~\bars{ejus \textit{streicht Hrsg.}} rursus centrum commune cum 4\textsuperscript{to} %
\textit{(2)}~id~\textit{L}}}
rursus cum centro quarti alicujus vel aliquot aliorum conjungatur, atque ita 
\edtext{tandem quotcunque}{\lemma{tandem}\Bfootnote{%
\textit{(1)}~omnium pa %
\textit{(2)}~quotcunque~\textit{L}}}
sphaerarum ac proinde cujuscunque corporis quod ex sphaeris interpolatis mente componi potest,\indexS{compositio corporis ex sphaeris interpolatis} aliquod erit centrum gravitatis.\indexS{centrum gravitatis} 
Possis etiam pro sphaera\indexS{sphaera} 
\edtext{adhibere sphaeroeidem\indexS{sphaeroeides}}{\lemma{}\Bfootnote{adhibere \bars{fusum Ellipticum vel \textit{gestr.}} sphaeroeidem~\textit{L}}}
vel etiam prisma\indexS{prisma} aliquod aut aliud ejusmodi 
\edtext{corpus, quod}{\lemma{corpus,}\Bfootnote{%
\textit{(1)}~in o %
\textit{(2)}~quod~\textit{L}}}
recta vel plano 
\edtext{per certum punctum}{\lemma{per}\Bfootnote{%
\textit{(1)}~aliquod 
\textit{(a)}~pu %
\textit{(b)}~cent %
\textit{(2)}~certum punctum~\textit{L}}}
transeunte semper in partes uniformes similes et aequales secatur.
\pend
\begin{edfigure}[1.5em]%
\hfill%
\hspace*{5mm}%
\includegraphics[scale=1]{%
gesamttex/edit_VIII-4/images/LH_37_05_021_d2_021v.pdf}%
\hfill%
\includegraphics[scale=1]{%
gesamttex/edit_VIII-4/images/LH_37_05_021_d3_021v.pdf}%
\hfill%
\edcaption{%
\lbrack\textit{Fig.~2}\rbrack%
\hspace*{41mm}%
\lbrack\textit{Fig.~3}\rbrack%
\hspace*{5mm}}
\end{edfigure}
